The event/trigger fabric (\peripheral{EVFAB}) is a PPI-style crossbar that lets peripherals command each other with no processor in the loop. Eight independent channels each hold one \{EVSEL, TASKSEL\} pair; when the event a channel selects fires and the channel is enabled, the fabric emits a registered one-\texttt{MCLK} pulse on the task it selects, one \texttt{MCLK} after the event. The whole block (the event front end, the crossbar, the output register, and the sticky status words) rides the free-running fabric clock (\texttt{MCLK}) in the always-on shared domain: \texttt{WFI} keeps it alive, field-power mode only slows it down, and \peripheral{PWRCTRL} never gates it. A chain therefore keeps firing with every hart asleep and the bus completely idle, which is the entire point of the block. \peripheral{EVFAB} is never a bus master, never stalls, and never rate-limits: it only ever adds one \texttt{MCLK} of latency between an event and the task pulse it drives. The main features are listed below:

\begin{itemize}
	\item Eight channels (\register{EVFCH0CFG}\ldots\register{EVFCH7CFG}), each a fixed \{EVSEL, TASKSEL\} pair; the channel-configuration array reserves sixteen addresses for future growth, so slots 8--15 read 0 and ignore writes on this build
	\item Sixteen wired event lines (\register{EVFEVSELn} encodes 5 bits; code 31 is the documented NONE), each carried into the fabric as a pulse, a toggle, or a level according to its source's clock domain; the fabric owns \textit{all} clock-domain crossing, so no producer peripheral needs its own synchronizer for this feature
	\item Ten wired task lines (\register{EVFTASKSELn} encodes 4 bits), delivered as clean registered one-\texttt{MCLK} pulses; a consumer never sees a held level from the fabric
	\item Write-1-to-set/write-1-to-clear channel-enable aliases (\register{EVFCHENSET}/\register{EVFCHENCLR}) so two harts can enable or disable different channels without racing a read-modify-write of \register{EVFCHEN}
	\item Sticky, write-1-to-clear \register{EVFFIRED} (channel fired), \register{EVFOVR} (a pulse was degraded), and \register{EVFEVSTAT} (a raw event was seen, recorded even while the fabric is disabled) status words, all set-wins over a same-cycle clear
	\item Firmware injection of a channel firing (\register{EVFCHTRIG}) or a raw event (\register{EVFEVTRIG}) with no producer hardware required at all, for bring-up and self-test
	\item Vectorless in this version: the interrupt output is a constant 0 and the \register{EVFIE} slot is reserved, so firmware polls \register{EVFSR}, whose two bits are live OR-reductions of \register{EVFFIRED} and \register{EVFOVR}
\end{itemize}

\subsection{Producers: the Event Lines}
\label{s:evfab-events}

Table \ref{t:evfab-events} lists every wired \register{EVFEVSELn} code. Every tap is taken from its source peripheral's flag-\textit{set} condition, before any interrupt mask is applied, so a chain works even with every relevant interrupt disabled at the source, and \register{EVFEVSTAT} can directly catch a tap that was mistakenly wired after a mask instead of before one. Each line carries one of three front-end modes, selected once per line at synthesis time and fixed for this build:

\begin{itemize}
	\item \textbf{Pulse (P):} the source is already a clean one-\texttt{MCLK} pulse in the \texttt{MCLK} domain, and the fabric passes it through combinationally (no flop, so the in-fabric latency for a pulse input is exactly the one output-register stage). A pulse wider than one \texttt{MCLK} cycle fires its channels once per cycle it is high, a producer contract, not a fabric limitation.
	\item \textbf{Toggle (T):} the source is a level that flips state in its own clock domain (a domain the fabric cannot assume is \texttt{MCLK}, or even free-running). The fabric two-flop synchronizes it into \texttt{MCLK} and XORs consecutive samples, producing exactly one \texttt{MCLK} pulse per flip in \textit{either} direction. A toggle input costs two extra \texttt{MCLK} cycles of latency versus a pulse input (three total from the flip to the task pulse).
	\item \textbf{Level (L):} the source is a level that the fabric two-flop synchronizes and edge-detects on the \textit{rising} edge only; a level that stays high forever fires exactly once, and the falling edge is silent.
\end{itemize}

\begin{table}[htb]
	\centering
	\begin{tabular}{ c l l l }
	\caption{\register{EVFEVSELn} event lines} \label{t:evfab-events}
	\textbf{EVSEL} & \textbf{Source} & \textbf{Mode} & \textbf{Notes} \\ \hline
	0  & \peripheral{RTC0} tick               & P & one-second periodic tick \\
	\rowcolor{tablehighlightcolor}
	1  & \peripheral{RTC0} alarm               & P & one-shot compare match \\
	2  & \peripheral{PWM0} period               & P & period-boundary tick \\
	\rowcolor{tablehighlightcolor}
	3  & \peripheral{PWM0} fault               & P & \register{FLTF} set condition \\
	4  & \peripheral{TIMER0} compare0           & T & \pin{T0CMP0} match toggle \\
	\rowcolor{tablehighlightcolor}
	5  & \peripheral{TIMER0} overflow           & T & counter-overflow toggle \\
	6  & \peripheral{TIMER1} compare0           & T & same shape as event 4, second instance \\
	\rowcolor{tablehighlightcolor}
	7  & \peripheral{UART0} receive             & P & receive-complete pulse \\
	8  & \peripheral{NFC0} field-detect         & T & carrier-presence toggle \\
	\rowcolor{tablehighlightcolor}
	9  & \peripheral{NFC0} rx-frame              & T & reader-frame-received toggle \\
	10 & \peripheral{DMA0} channel-0 done        & P & \register{DMADONE[0]} set pulse \\
	\rowcolor{tablehighlightcolor}
	11 & \peripheral{DMA0} channel-1 done        & P & \register{DMADONE[1]} set pulse \\
	12 & \peripheral{DMA0} error (combined)      & P & OR of all channel error-set pulses \\
	\rowcolor{tablehighlightcolor}
	13 & \peripheral{TRNG0} data-ready           & L & \register{TRNGDRDY} level \\
	14 & \peripheral{I2CT0} address-match        & P & \register{I2CTAMF} set pulse \\
	\rowcolor{tablehighlightcolor}
	15 & \peripheral{GPIO0} masked pad edge      & (special) & see below \\
	16--31 & reserved                           & -- & always inert (no line drives them) \\
	\hline
	\end{tabular}
\end{table}

\noindent Event 15 does not come from a plain input line: it is built \textit{inside} the fabric from \peripheral{GPIO0}'s eight raw pad levels (the pre-mask value, independent of every pin's \register{PxIE} interrupt enable). Each of the eight pads is independently two-flop synchronized and rising-edge detected, ANDed against its own bit in \register{EVFGPIOMASK}, and the eight results are OR-reduced onto event 15; a bit clear in \register{EVFGPIOMASK} means that pad's edges never reach the fabric at all, which also means a masked-off, floating, or unbonded pad can never inject an unknown value into the crossbar. \register{EVFGPIOMASK} resets to 0, so this path is completely inert out of reset.

Events whose source peripheral is a config-droppable block are only wired in configurations that include that block: \peripheral{RTC0} (event 0/1), \peripheral{PWM0} (event 2/3), \peripheral{TIMER1} (event 6), \peripheral{NFC0} (event 8/9), \peripheral{DMA0} (event 10/11/12), \peripheral{TRNG0} (event 13), and \peripheral{I2CT0} (event 14). In a configuration that drops one of these blocks its event line is tied inactive, so a channel programmed with that \register{EVFEVSELn} code is simply and permanently inert, exactly as if it had been programmed with the reserved NONE code. \peripheral{TIMER0}, \peripheral{UART0}, and \peripheral{GPIO0} are always present, so events 4, 5, 7, and 15 are wired in every configuration.

The \register{EVFEVSELn} codes are an ABI: once a producer ships, its code number never moves, and a firmware image that only knows the codes it uses is unaffected by every future addition to this table. Codes 16 through 31 are reserved for future producers (deferred \peripheral{TIMER0}/\peripheral{TIMER1} input captures, the remaining \peripheral{GPIOx} ports, \peripheral{QSPI0}, \peripheral{OW0}, \peripheral{I3C0}, \peripheral{UART0}/\peripheral{UART1} transmit events, \peripheral{NFC0} transmit-done/CRC-error, further \peripheral{DMA0} channels, and a \peripheral{TRNG0} health alarm); programming one of them, or code 31 (the documented NONE encoding), leaves a channel structurally inert with no special-case behavior anywhere in the fabric.

\subsection{Consumers: the Task Lines}
\label{s:evfab-tasks}

Table \ref{t:evfab-tasks} lists every wired \register{EVFTASKSELn} code. Every task input is, by construction, a clean registered one-\texttt{MCLK} pulse: the single output-register stage described in Section \ref{s:evfab-map} is the \textit{only} flop between an event and every task line, so a consumer can never observe a held level from the fabric and never needs its own pulse-stretcher. Where a task's own action would be hazardous if delivered as a held level (for example, a toggle that would produce a continuous square wave rather than one event), that shape is simply not offered as a task in this version.

\begin{table}[htb]
	\centering
	\begin{tabular}{ c l l }
	\caption{\register{EVFTASKSELn} task lines} \label{t:evfab-tasks}
	\textbf{TASKSEL} & \textbf{Consumer} & \textbf{Action} \\ \hline
	0  & \peripheral{DMA0} channel-0 GO   & launches channel 0 exactly like a \register{DMAGO} write \\
	\rowcolor{tablehighlightcolor}
	1  & \peripheral{DMA0} channel-1 GO   & launches channel 1 exactly like a \register{DMAGO} write \\
	2  & \peripheral{TIMER0} START        & sets the timer's run bit \\
	\rowcolor{tablehighlightcolor}
	3  & \peripheral{TIMER0} STOP         & clears the timer's run bit \\
	4  & \peripheral{PWM0} fault trip     & equivalent to a \register{FLTTRIG} write \\
	\rowcolor{tablehighlightcolor}
	5  & \peripheral{PWRCTRL} tile wake   & clears the addressed tile's held-in-gate request \\
	6  & \peripheral{NPU0} THINK          & starts a computation exactly like an \register{NPUCR} bit-16 write \\
	\rowcolor{tablehighlightcolor}
	7  & \peripheral{GPIO0} OUT-SET       & sets every pin selected in \register{PxTASK} \\
	8  & \peripheral{GPIO0} OUT-CLR       & clears every pin selected in \register{PxTASK} \\
	\rowcolor{tablehighlightcolor}
	9--15 & reserved                     & \register{EVFFIRED} still sets; no task line exists \\
	\hline
	\end{tabular}
\end{table}

\noindent Tasks 0/1 (\peripheral{DMA0}), 4 (\peripheral{PWM0}), and 6 (\peripheral{NPU0}) are only wired in configurations that include those blocks; \peripheral{TIMER0}, \peripheral{PWRCTRL}, and \peripheral{GPIO0} are always present, so tasks 2, 3, 5, 7, and 8 are wired in every configuration. A channel programmed with \register{EVFTASKSELn} 9 through 15, or targeting a task whose consumer block is absent from the configuration, still sets \register{EVFFIRED} on a matching event (the fabric decoded a real firing) but drives no task line and can never set \register{EVFOVR} for that channel; there is no consumer busy line to compare against, so the fabric correctly reports nothing was degraded.

A firing channel launches its task action \textit{in addition to}, not instead of, the consumer's ordinary software-driven path: a task-driven \register{DMA0} GO behaves exactly like a register \register{DMAGO} write (subject to the same busy/enable checks), a task-driven \peripheral{TIMER0} START or STOP behaves exactly like the corresponding control-register write, and so on. Where a task and a same-cycle CPU write could disagree on the same bit, the peripheral's own chapter states which one wins (\peripheral{GPIO0}'s \register{PxTASK} outputs, for example, apply after the CPU's write, and a same-cycle SET+CLEAR on the same pin resolves to CLEAR).

\textbf{No task consumer counts firings.} Two or more channels landing on the same task in the same cycle merge into exactly one task pulse (recorded as \register{EVFOVR} on every colliding channel, Section \ref{s:evfab-map}); every wired task line is safe to receive that merged single pulse, because none of them has counting semantics: a GO, a START, a STOP, a fault trip, a tile-wake, a THINK, or a pin set/clear are all idempotent-per-cycle actions. This is a binding property of the six touched blocks, not something a driver must additionally guard: if a future task addition ever needed to count distinct firings rather than react to ``at least one fired,'' it could not use this fabric's task interface as-is.

\subsection{Channel Enable, Status, and Injection}
\label{s:evfab-map}

The fabric lives at \texttt{0x6B00} (Section \ref{s:evfab-shared}), decoded as sixty-four consecutive words. \register{EVFCR} bit 0 is the single global enable; \register{EVFCHEN} then gates each of the eight channels independently. A channel fires only when \register{EVFCR.EN}, its own \register{EVFCHEN} bit, and its selected event are all true; clearing either gate makes the channel \textit{completely} inert: no task pulse, no \register{EVFFIRED}, no \register{EVFOVR}. \register{EVFEVSTAT} sits upstream of both gates and keeps recording regardless, which is deliberate: it lets firmware (or a bench) prove that an event tap was taken \textit{before} any mask, by disabling everything downstream and still observing \register{EVFEVSTAT} set.

\register{EVFCHEN} may be written directly, but firmware normally uses the write-1-to-set/write-1-to-clear aliases \register{EVFCHENSET} and \register{EVFCHENCLR} instead: each is a mask where writing a 1 to bit $n$ sets or clears \register{EVFCHEN}$[n]$ and writing 0 leaves that bit untouched, and both read back the live \register{EVFCHEN} value. Two harts enabling or disabling \textit{different} channels can therefore each issue one write with no read-modify-write race and no risk of one hart's write clobbering the other's in-flight change, the same convention as \register{PxOUTS}/\register{PxOUTC} and the \peripheral{IRQROUTER} claim mechanism. Each \register{EVFCHnCFG} word carries that channel's \register{EVFEVSELn} field, its \register{EVFTASKSELn} field, and a read-only \register{EVFENRn} mirror of its live \register{EVFCHEN} bit, so a driver can read a channel's full configuration and enable state back in one access.

\register{EVFFIRED} (one sticky bit per channel) sets whenever an enabled channel fires and clears only on an explicit write-1; a firing that lands in the same cycle as a write-1-to-clear leaves the bit \textit{set} (set-wins), so a status word being drained can never silently swallow an event that arrives mid-clear. \register{EVFOVR} sets, per firing channel, when that channel's task consumer was already busy at the moment of firing, or when two or more enabled channels landed on the same task in the same cycle (Section \ref{s:evfab-tasks}). In every such case the task pulse is still emitted normally; \register{EVFOVR} only records that the pulse was \textit{degraded} in some way, never that it was withheld. Both are drained the same way: read to see which channels are set, write 1 to the bits just handled.

\register{EVFCHTRIG} and \register{EVFEVTRIG} are write-only injection registers (both read 0) that let firmware exercise the fabric with no producer hardware involved at all: a 1 written to bit $n$ of \register{EVFCHTRIG} injects one firing directly at channel $n$, honoring that channel's \register{EVFCR}/\register{EVFCHEN} gates exactly like a real event (a CHTRIG to a disabled channel does nothing); a 1 written to bit $e$ of \register{EVFEVTRIG} injects one occurrence of raw event $e$, which is recorded in \register{EVFEVSTAT} even with the fabric globally disabled and reaches the crossbar exactly like a genuine tap. Both share one hardware property that also protects real bus writes: the access itself may be held on the bus for several clock edges (a slow or contended master), but the fabric decodes the write \textit{once} per select window, so a single CPU store, however long the underlying bus access takes to retire, produces exactly one injection, never one per bus-clock edge it happens to span.

\register{EVFCAP} is a read-only word that reports the \textit{live} counts this build implements (8 channels, 16 events, 10 tasks, fabric version 1), so a driver can size its loops from the register instead of a compile-time constant, and instead of the wider \textit{encode spaces} that \register{EVFEVSELn}/\register{EVFTASKSELn} reserve for growth.

\subsection{Driver Contract}
\label{s:evfab-contract}

\begin{itemize}
	\item \textbf{Configure before enabling.} \register{EVFCHnCFG}.\register{EVFEVSELn} and \register{EVFCHnCFG}.\register{EVFTASKSELn} must only be changed while that channel's \register{EVFCHEN} bit is 0; the fabric does not interlock a live reconfiguration against an in-flight firing, so the sequence is: clear the channel (\register{EVFCHENCLR}), write its new \register{EVFCHnCFG}, then re-enable it (\register{EVFCHENSET}).
	\item \textbf{Action-visibility latency.} A write to \register{EVFCHTRIG}, \register{EVFEVTRIG}, or any of the write-1-to-clear status words (\register{EVFFIRED}, \register{EVFOVR}, \register{EVFEVSTAT}) takes effect three \texttt{MCLK} cycles after its bus access opens; the register file itself is on the gated bus clock, but these particular actions are non-idempotent and are decoded once, in the free-running \texttt{MCLK} domain, by a dedicated edge detector. A status read issued immediately after such a write (only possible from a master faster than the shared multi-core bus) can observe stale state; on \AsicNameForUserGuide{} a shared-window access already costs several \texttt{MCLK} cycles end to end, so this window is not observable in practice, and no special handling is needed beyond not assuming same-instruction visibility.
	\item \textbf{No consumer may count firings.} Section \ref{s:evfab-tasks}'s merge rule means a driver must not build a protocol on top of a task line that depends on receiving one pulse per channel firing when two channels share that task. Design channel-to-task assignments so every task is either driven by one channel, or every colliding channel's action is genuinely idempotent (the case for all ten wired tasks today).
	\item \textbf{The NPU staging-RAM contract still applies to task 6.} A fabric-driven \peripheral{NPU0} \texttt{THINK} obeys exactly the same rule as a register-driven one: no hart may touch \texttt{0xC000}--\texttt{0xFFFF} while a \texttt{THINK} may be in flight. Because the fabric can start a \texttt{THINK} with no CPU instruction executing at all, a design that uses task 6 must guarantee by construction, not by polling right before the event, that nothing else is touching the staging RAM when the triggering event might occur; there is no hardware interlock.
	\item \textbf{Channel ownership across harts is a software convention.} There is no \register{OWNER} field anywhere in the register map. Multiple harts sharing the fabric should partition the eight channels by index (for example, by convention rather than by any register), using \register{EVFCHENSET}/\register{EVFCHENCLR} so that enabling or disabling one hart's channels never touches another's, and never writing another hart's \register{EVFCHnCFG} without an out-of-band handshake (a mutex, or simply "this channel belongs to hart 2").
\end{itemize}

\subsection{Worked Example: An Autonomous Sensing Chain}
\label{s:evfab-example}

The following configuration lets \peripheral{TIMER0} trigger a \peripheral{UART0} transmission through \peripheral{DMA0} with the CPU asleep the entire time, and is the reference chain proven by the \texttt{shevfab} test in the behavioral regression:

\begin{center}
\texttt{TIMER0 compare0 (EV4)} $\rightarrow$ \texttt{channel 0} $\rightarrow$ \texttt{DMA0 channel-0 GO (TASK 0)} $\rightarrow$ \texttt{DMA0 fills UART0TX through the arbiter} $\rightarrow$ \texttt{DMA0 channel-0 done (vector 118)} $\rightarrow$ \texttt{IRQROUTER} $\rightarrow$ \texttt{meip} $\rightarrow$ \texttt{hart wakes}
\end{center}

Firmware programs \register{EVFCH0CFG} with \register{EVFEVSELn} = 4 (\peripheral{TIMER0} compare0) and \register{EVFTASKSELn} = 0 (\peripheral{DMA0} channel-0 GO) while channel 0 is disabled, arms \peripheral{DMA0} channel 0 with its source/destination/length pointed at the \peripheral{UART0} transmit path and \textit{its own} \register{DMAGO} bit left clear (so the fabric task pulse is the \textit{only} way the transfer launches; a broken event path cannot be masked by an accidental register GO left over from a previous test), sets \register{EVFCR.EN} and \register{EVFCHENSET} bit 0 (arming \peripheral{TIMER0}'s own compare0 interrupt path is not required at all; the fabric taps the compare event directly, before any interrupt enable), enables \peripheral{DMA0}'s done interrupt, routes vector 118 to itself in \peripheral{IRQROUTER}, starts \peripheral{TIMER0}, and executes \texttt{WFI}.

From that point the hart executes nothing until the chain completes: \peripheral{TIMER0}'s compare0 toggle reaches the fabric three \texttt{MCLK} cycles after it flips (Section \ref{s:evfab-events}, toggle mode), fires channel 0, and one \texttt{MCLK} later the fabric drives a task pulse into \peripheral{DMA0}'s channel-0 GO input exactly as if firmware had written \register{DMAGO} itself. \peripheral{DMA0} runs its transfer over the shared arbiter, sets its channel-0 done flag on completion, the combined done interrupt reaches the router, the router raises \texttt{meip} to the hart, and the hart's interrupt-service routine clears \register{DMADONE}\texttt{[0]} and resumes normal execution. Every link in that path (the event tap, the fabric crossbar, the task action, the DMA transfer, and the interrupt delivery) is independently a documented, testable contract; the chain is simply those contracts composed with nobody polling in between.

A companion pattern from the same design worth noting, though it is not shipped as a single test: \peripheral{RTC0}'s periodic tick (event 0) driving \peripheral{TIMER0} START (task 2), with a second channel wiring \peripheral{TIMER0}'s own compare event back to \peripheral{TIMER0} STOP (task 3), produces an autonomous, CPU-free sensing window (the real-time clock opens a timed measurement interval and the timer itself closes it) entirely inside the fabric.

\subsection{Shared Access on \AsicNameForUserGuide}
\label{s:evfab-shared}

In configurations that include it, \peripheral{EVFAB} lives in the shared peripheral window behind the multi-core arbiter (see Section \ref{s:multicore}) in the mutex-bank page, sub-slot 11 at \texttt{0x6B00}, so any hart can configure it and any hart's channels can drive any wired task. \peripheral{EVFAB} has no pins and, because its datapath and register file both ultimately ride the free-running fabric clock, its operation is immune to clock reconfiguration: software does \emph{not} need to write \register{\register{SYSCLKCR}} to 0 before using it, unlike the \texttt{SMCLK}-domain serial peripherals. This version spends no interrupt vector at all: \texttt{irq\_evfab} is wired to a constant 0 and the \register{EVFIE} slot is reserved and reads 0, so \peripheral{EVFAB} never appears in the \peripheral{IRQROUTER}'s vector list and the frozen interrupt-vector numbering is completely undisturbed by this peripheral. A future version that spends a vector for \register{EVFFIRED}/\register{EVFOVR} would be a purely additive change (implement \register{EVFIE}, drive \texttt{irq\_evfab}, and route the new vector) with no register-map break.
